\documentclass{ximera}
%geschreven 19/11/26
\title{Herhalingsoefeningen elektrische krachtwerking, veld en potentiaal}
\author{Daan Lipkens}

\begin{document}
\begin{abstract}
    Herhalingsoefeningen over hoofdstuk 1 tot en met 3, als voorbereiding op de syntheseproef.
\end{abstract}
\maketitle

In dit deel herhaal je de leerstof van de elektrische krachtwerking, het elektrisch veld en de potentiaal. In het eerste onderdeel staan vragen en vraagstukken door elkaar: bij elke oefening moet je zelf herkennen over welke leerstof het gaat en welke aanpak past. In het tweede onderdeel staan vraagstukken die leerstof uit meerdere hoofdstukken combineren. Maak de oefeningen eerst zonder hints en zonder de cursus, zoals op de syntheseproef. De coulombconstante $k$ en de elementaire lading $e$ ken je van buiten. Andere gegevens, zoals de massa van een deeltje, staan in de opgave.

\end{document}
